\section*{Your turn}

Check out this chapter's tag and confirm you have the same service I do:

\begin{minted}{text}
git clone https://github.com/Giang-Dang/mosaic-graph
cd mosaic-graph
git checkout ch03
pwsh scripts/verify.ps1
\end{minted}

The script now checks two more things than it did in
chapter~\ref{ch:hotchocolate-quickly}: that the pipeline is those thirteen
middleware in that order, and that the timeline reports 146 resolvers. Run the
service with its console where you can see it for all of these.

\begin{enumerate}
  \item \textbf{Watch a document go cold.} Send the catalog-with-reviews query
  three times and read the three timeline lines. Now change one character of it,
  a field alias will do, and send it again. Which phases came back, and which
  stayed dashes? The answer tells you what the cache is keyed on.

  \item \textbf{Find the request that leaves no trace.} Post a document with an
  unclosed brace, then post one asking for a field that does not exist. Both
  come back 400. Count the timeline lines. Then explain the difference to
  somebody without using the word \enquote{validation}.

  \item \textbf{Measure the caches honestly.} Restart the service and time the
  first request. Then send twenty, and time the twenty-first. Then send a novel
  document and time that. You now have three numbers, only two of which belong
  in the same sentence. Decide which two, and write down why the third is
  misleading.

  \item \textbf{Insert your own middleware.} Add one that logs the operation
  name, keyed, and anchored \code{after:} the document validation middleware.
  Confirm it appears in the startup listing where you expected. Then move its
  registration above \code{AddGraphQL()} and see what the listing says. This is
  the mistake section~\ref{sec:ch03-pipeline} describes, and it is cheaper to
  make on purpose.

  \item \textbf{Break the listener.} Delete
  \code{AddApplicationService<ILoggerFactory>()} from \code{Program.cs} and
  start the service. Read the error. Then work out, without running it, whether
  the same mistake in an error filter would fail the same way.

  \item \textbf{Move a field's scope.} Run the \code{samples/resolver-scopes}
  project and send the four-field query from
  section~\ref{sec:ch03-scopes}. Predict the numbers before you look. Then add
  \code{[UseResolverScope]} to the field that currently has
  \code{[UseRequestScope]}, and predict them again.

  \item \textbf{Cost something.} Send three queries of increasing size with the
  \code{GraphQL-Cost: validate} header and record the field costs. Work out the
  rule from the numbers, then check it against the \code{@cost} weights in
  \code{schema/mosaic.graphql}. Which fields are free, and is that the answer
  you would want if you were setting a budget?
\end{enumerate}

Exercise~2 is the one that matters at three in the morning. A request that
returns 400 without producing a single log line, span or diagnostic event is not
a broken server, and the instinct that it must be will send you looking in the
executor for something the transport rejected. From
chapter~\ref{ch:enter-the-router} onward there is another process in front of
these services, with its own layers and its own version of the same trap, and
knowing which of your components can answer without recording anything is most
of the skill of reading a distributed failure.
